\section{复数的几何意义}

\begin{definition}[复数的几何意义]
    复数 $z=a+bi(a,b\in \mathbb{R})$ 与复平面内的点 $Z(a,b)$ 一一对应, 与复平面内的向量 $\overrightarrow{OZ}$ 一一对应。
\end{definition}

\begin{example}
    已知 $m \in \mathbf{R}, \mathrm{i}$ 是虚数单位，复数 $z=m^2+m-2+\left(m^2-1\right) \mathrm{i}$ ．
    \begin{enumerate}
        \item 若复数 $z$ 对应的点位于第二象限，求 $m$ 的取值范围．
    \end{enumerate}
\end{example}

\v{8em}

\begin{example}
    在复平面内, 复数 $\frac{a+2i}{i}(a\in \mathbb{R})$ 对应的点在第四象限, 则实数a的取值范围是( \quad )
    \begin{tasks}(4)
        \task $(0,+\infty)$
        \task $(-\infty,0)$
        \task $(2,+\infty)$
        \task $(-\infty,2)$
    \end{tasks}
\end{example}

\v{4em}

\begin{example}
    已知复数 $z=\frac{1+\sqrt{3} i}{\sqrt{3}+i}$（其中 $i$ 为虚数单位），则 $z$ 在复平面内对应的点位于 （\qquad）
    \begin{tasks}(4)
        \task 第一象限
        \task 第二象限
        \task 第三象限
        \task 第四象限
    \end{tasks}
\end{example}

\v{4em}

\begin{example}
    已知复数 $z=2+\frac{1+i}{1-i}$（ $i$ 为虚数单位），则 （\qquad）
    \begin{tasks}(2)
        \task $z$ 的共轭复数 $z$ 的虚部为 $-i$
        \task $z-2$ 为纯虚数
        \task $z^2$ 的模为 5
        \task! 若在复平面内，向量 $\overrightarrow{A B}$ 对应的复数为 $z$ ，向量 $\overrightarrow{C B}$ 对应的复数为 $-1+3 i$ ，则向量 $\overrightarrow{C A}$ 对应的复数为 $-3+2 i$
    \end{tasks}
\end{example}

\v{8em}

\subsection{复数和向量结合}

\begin{example}
    在复平面内, O为坐标原点, 复数 $z_{1}=i(-4+3i)$, $z_{2}=7+i$ 对应的点分别为 $Z_{1}$, $Z_{2}$, 则 $\angle Z_{1}OZ_{2}$ 的大小为( \quad )
    \begin{tasks}(4)
        \task $\frac{\pi}{3}$
        \task $\frac{2\pi}{3}$
        \task $\frac{3\pi}{4}$
        \task $\frac{5\pi}{6}$
    \end{tasks}
\end{example}

\vspace{4em}

\section{复数的模长}

\begin{definition}[复数的模]
    设复数 $z=a+bi(a,b\in \mathbb{R})$, 我们把向量 $\overrightarrow{OZ}$ 的模叫做z的模, 记作 $|z|$, 且 $|z|=\sqrt{a^{2}+b^{2}}$.
\end{definition}

\begin{example}
    设 $i$ 为虚数单位, 已知复数 $z=\frac{5}{a+i}$ 满足 $|z|=\sqrt{5}$, 其中 $a\in \mathbb{R}$ 且 $a>0$, 则 $z= \underline{\hspace{2cm}}$.
\end{example}

\vspace{4em}

\begin{example}
    已知 $i$ 为虚数单位, 复数满足 $|z-2i|=|z|$, 则 $z$ 的虚部为 \underline{\hspace{2cm}}.
\end{example}

\v{4em}

\begin{example}(2022 全国乙卷)
    已知 $z=1-2i$, 且 $z+a\overline{z}+b=0$, 其中a,b为实数, 则( \quad )
    \begin{tasks}(4)
        \task $a=1$, $b=-2$
        \task $a=-1$, $b=2$
        \task $a=1$, $b=2$
        \task $a=-1$, $b=-2$
    \end{tasks}
\end{example}

\vspace{4em}

\subsection{复数的模有关的推论}

\subsubsection{共轭模长推论}

\begin{theorem}[共轭模长推论]
    \begin{itemize}
        \item $|z|^{2}\ne z^{2}$. (设 $z=a+bi(a,b\in \mathbb{R},b\ne0)$, 则 $|z|^{2}=a^{2}+b^{2}$, $z^{2}=a^{2}-b^{2}+2abi$, 显然不等)
        \item  $z\cdot\overline{z}=|z|^{2}=|\overline{z}|^{2}$ 是共轭复数的重要性质。
    \end{itemize}

    \begin{tips}
        设 $z=a+bi$, 则 $\overline{z}=a-bi$, 所以 $|z|=|\overline{z}|=\sqrt{a^{2}+b^{2}}$, $z\cdot\overline{z}=(a+bi)(a-bi)=a^{2}-b^{2}i^{2}=a^{2}+b^{2}$
    \end{tips}
\end{theorem}

\begin{example}
    已知$\overline{z}$是复数$z$的共轭复数, 则下列式子中与$z\cdot\overline{z}$不相等的是( \quad )
    \begin{tasks}(4)
        \task $|\overline{z}^{2}|$ \\
        \task $|z|^{2}$ \\
        \task $|z^{2}|$ \\
        \task $\overline{z}^{2}$
    \end{tasks}
\end{example}

\begin{solution}
    动用上述结论可以干掉前两个选项，当然我们也可以严格设数来证明。
    \begin{itemize}
        \item 对于 A 选项，$|\overline{z}^{2}|=|\overline{z}\cdot\overline{z}|=|\overline{z}|\cdot|\overline{z}|=|\overline{z}|^{2}=a^{2}+b^{2}=z\cdot\overline{z}$;
        \item 对于 B 选项，因为 $|z|^{2}=a^{2}+b^{2}$, 结合（1）知 $z\cdot\overline{z}=|z|^{2}$;
        \item 对于 C 选项，$|z^{2}|=|z\cdot z|=|z|\cdot|z|=|z|^{2}=a^{2}+b^{2}$, 结合（1）知 $|z^{2}|=z\cdot\overline{z}$;
        \item 对于 D 选项，$\overline{z}^{2}=(a-bi)^{2}=a^{2}+b^{2}i^{2}-2abi=a^{2}-b^{2}-2abi\ne z\cdot\overline{z}$, 故选D.
    \end{itemize}
\end{solution}

\vspace{4em}

\begin{example}(2022 全国甲卷)
    若 $z=-1+\sqrt{3}i$, 则 $\frac{z}{z\overline{z}-1}=($ \quad $)$
    \begin{tasks}(4)
        \task $-1+\sqrt{3}i$
        \task $-1-\sqrt{3}i$
        \task $-\frac{1}{3}+\frac{\sqrt{3}}{3}i$
        \task $-\frac{1}{3}-\frac{\sqrt{3}}{3}i$
    \end{tasks}
\end{example}

\vspace{4em}

\begin{example}(多选)
    设 $z_{1}$, $z_{2}$, $z_{3}$ 为复数, $z_{1}\ne0$, 下列命题中正确的是( \quad )
    \begin{itemize}
        \item[A.] 若 $|z_{2}|=|z_{3}|$, 则 $z_{2}=\pm z_{3}$
        \item[B.] 若 $z_{1}z_{2}=z_{1}z_{3}$, 则 $z_{2}=z_{3}$
        \item[C.] 若 $\overline{z}_{2}=z_{3}$, 则 $|z_{1}z_{2}|=|z_{1}z_{3}|$
        \item[D.] 若 $z_{1}z_{2}=|z_{1}|^{2}$, 则 $z_{1}=z_{2}$
    \end{itemize}
    \begin{solution}
        \textcolor{blue}{对于 C 选项，也可以通过纯代数的形式求解。}
        \begin{itemize}
            \item 对于 A，$|z_{2}|=|z_{3}|$ 可看成复平面内 $|\overline{OZ_{2}}|=|\overline{OZ_{3}}|$, 但方向未定, 故 $z_2 = \pm z_3$ 不一定成立。举个反例, 取 $z_{2}=\sqrt{3}+i$, $z_{3}=1+\sqrt{3}i$, 则 $|z_{2}|=|z_{3}|=2$, 但 $z_{2}\ne\pm z_{3}$, 故A项错误;
            \item 对于 B，由 $z_{1}z_{2}=z_{1}z_{3}$ 两端同除以非零复数 $z_1$ 可得 $z_{2}=z_{3}$, 故B项正确;
            \item 对于 C，看到共轭, 想到模的性质。因为 $\overline{z}_{2}=z_{3}$, 所以 $|\overline{z}_{2}|=|z_{3}|$, 又 $|\overline{z}_{2}|=|z_{2}|$, 所以 $|z_{2}|=|z_{3}|$, 从而 $|z_{1}z_{2}|-|z_{1}z_{3}|=|z_{1}|\cdot|z_{2}|-|z_{1}|\cdot|z_{3}|=|z_{1}|(|z_{2}|-|z_{3}|)=0$, 故 $|z_{1}z_{2}|=|z_{1}z_{3}|$, 故C项正确;
            \item 对于 D，我们知道, $z_{1}\cdot\overline{z}_{1}=|z_{1}|^{2}$, 故要使 $z_1z_2 = |z_1|^2$, 只需 $z_2=\overline{z}_1$, 而不是 $z_{2}=z_{1}$。下面举个反例, 取 $z_{1}=1+i$, $z_{2}=1-i$, 满足 $z_{1}z_{2}=(1+i)(1-i)=1-i^{2}=2=|z_{1}|^{2}$, 但 $z_{1}\ne z_{2}$, 故D项错误。
        \end{itemize}
        \begin{tips}
            对于C，设 $z_{1}=a+bi$, $z_{2}=c+di$, 其中 $a,b,c,d\in \mathbb{R}$, 因为 $\overline{z}_{2}=z_{3}$, 所以 $z_{3}=c-di$。
            $$
                \begin{aligned}
                    |z_{1}z_{2}|=|(a+bi)(c+di)|=|(ac-bd)+(ad+bc)i|=\sqrt{(ac-bd)^{2}+(ad+bc)^{2}} \\
                    =\sqrt{a^{2}c^{2}-2acbd+b^{2}d^{2}+a^{2}d^{2}+2adbc+b^{2}c^{2}}
                \end{aligned}
            $$

            $$
                \begin{aligned}
                    |z_{1}z_{3}|=|(a+bi)(c-di)|=|(ac+bd)+(bc-ad)i|=\sqrt{(ac+bd)^{2}+(bc-ad)^{2}} \\
                    =\sqrt{a^{2}c^{2}+2acbd+b^{2}d^{2}+b^{2}c^{2}-2bcad+a^{2}d^{2}}
                \end{aligned}
            $$
        \end{tips}
    \end{solution}
\end{example}

\begin{example}(2023 浙江温州三模)(多选)
    已知复数 $z_1, z_2$, 下列命题正确的是( \quad )
    \begin{tasks}(4)
        \task $|z_{1}z_{2}|=|z_{1}|\cdot|z_{2}|$
        \task 若 $|z_{1}|=|z_{2}|$, 则 $z_{1}=\pm z_{2}$
        \task $z_{1}\overline{z_{1}}=|z_{1}|^{2}$
        \task 若 $z_{1}^{2}=\overline{z}_{1}^{2}$, 则 $z_1$ 为实数
    \end{tasks}
\end{example}

\v{6em}

\subsubsection{乘除模长推论}

\begin{theorem}[乘除模长推论]
    已知 $z_1=a+b \mathrm{i}(a, b \in \mathrm{R}), z_2=c+d \mathrm{i}(c, d \in \mathrm{R})$ 且 $\left(c^2+d^2 \neq 0\right)$。则：
    $$\left|\frac{z_1}{z_2}\right|=\frac{\left|z_1\right|}{\left|z_2\right|}=\frac{\sqrt{a^2+b^2}}{\sqrt{c^2+d^2}}\left(c^2+d^2 \neq 0\right),\left|z_1 \cdot z_2\right|=\left|z_1\right| \cdot\left|z_2\right|=\sqrt{a^2+b^2} \cdot \sqrt{c^2+d^2}$$
    根据上式可以推出：
    $$|z|^n=|z^n|=|z| \cdot|z| \cdots|z|=|z|^n(n \in \mathrm{N}^*)$$
\end{theorem}

\begin{example}
    若复数 $z=(1-i)(1+2 i)$ ，则复数 $z$ 的虚部为 \underline{\qquad}; $\left|\frac{(1-i)(1+2 i)}{1+\sqrt{3} i}\right|=$ \underline{\qquad}.
\end{example}

\v{8em}

\begin{example}
    (15年高二下徐汇期中)
    若 $z=\frac{(1+i)^4(-1-\sqrt{3} i)^7}{(1-i)^{12}}$, 则 $|z|=$ \underline{\qquad}.
\end{example}

\v{8em}
